\cotefile{156-9}
\scanpage{105}
\mspage{\unsure{151|101}}
% archive sheet number: 104

Ces questions ne \unsure{paraissent} pas, à $1^{\mathrm{ère}}$ vue du moins, être \unsure{de|des} \unsure{conséquence|conséquences} \unsure{directe|directes|d'intérêt} pour le \unsure{formalisme} ensembliste des ateliers. Il n'en va plus de même des suivantes.

\uline{Analogie ensembliste}

Soient $E$, $E'$ des ensembles, et
\[
\unsure{\overleftarrow{\alpha}_*|\widetilde{\alpha}_*}:
\unsure{\mathcal{P}|\mathscr{P}}(E)\mathrel{\unsure{-|\longrightarrow}}\mathcal{P}(E')
\]
une application associée à une corr. $\alpha$, i.e. $\overleftarrow{\alpha}_*$ commute aux Sup \unsure{quelc.} Soit $\Gamma$ le graphe, \unsure{avec}
\diagram{\begin{tikzcd}[ampersand replacement=\&]
\Gamma \arrow[r,"f'"'] \arrow[d,"f"] \& E' \\
E \&
\end{tikzcd}}
\unsure{d'où} $\bar\alpha(\deleted{\unsure{E}}A)=f'(f^{-1}(A))$.

\textcircled{A} Conditions équivalentes
\begin{enumerate}[label={},leftmargin=1.8em]
\item (i) $\overleftarrow{\alpha}_*$ commute à $|\unsure{\circ|0|\bullet}|$
\item (ii) $\overleftarrow{\alpha}_*$ commute à $S\wedge T\quad \operatorname{Inf}_{i\in I}S_i$, \deleted{pour} $I\ne\varnothing$
\item \inserted{(iii) $\alpha_*$ \unsure{commute} :}
\item (iii') $\overleftarrow{\alpha}_*$ commute à $S\wedge\unsure{C|\complement}T$
\item (iii'') $\overleftarrow{\alpha}_*$ commute : $S\wedge CT$ pour $T\leq S$
\item (iv) $f'$ injection --- i.e. $\bar\alpha$ \unsure{composé} d'un comorphisme, et d'un ``morphisme'' associé à une application injection (ou ce qui revient au même, définissant une application $\bar f'$ injection $\mathcal{P}(\Gamma)\mathrel{-}\mathcal{P}(E')$.
\end{enumerate}

\uline{Corollaire 1} Conditions équivalentes
% margin note written sideways
% slanted
\margin{\unsure{découle} \unsure{après} $B$ cor. \unsure{de} $A$ et $B$ \unsure{$\mathrm{i)}\longrightarrow$|$\mathrm{ii)}\longrightarrow$}}
\begin{enumerate}[label={},leftmargin=1.8em]
\item (i) $\overleftarrow{\alpha}_*$ commute à $|\circ|$ et $\overleftarrow{\alpha}_*(\unsure{1|A})=1$
\item (ii) $\overleftarrow{\alpha}_*\quad\rule{3.8cm}{0.3pt}\quad S\wedge T$ et $\overleftarrow{\alpha}_*(1)=1$
\item (\unsure{ii'|iii}) $\alpha_*\quad\rule{3cm}{0.3pt}\quad\operatorname{Inf}_{i\in I}S_i$, $I$ fini
\end{enumerate}

\scanpage{106}
\mspage{102}
% archive sheet number: 105

\begin{enumerate}[label={},leftmargin=1.8em]
\item \unsure{(ii'')} $\unsure{\overleftarrow{\alpha}_*|\widetilde{\alpha}_*}$ commute aux $\operatorname{Inf}$ quelconques
\item (iii) $\overleftarrow{\alpha}_*$ commute à $\unsure{C|\complement}$
\item (iv) $\alpha$ est un comorphisme, \inserted{i.e. $f':\Gamma\unsure{\simeq|\mathrel{-}}E'$} (i.e. $\unsure{\alpha_*|\bar\alpha_*}$ de la forme $A\mathrel{\unsure{-|\longmapsto}}g^{-1}(A)$, où $g\,(=f f'^{-1}):E'\mathrel{\unsure{-|\longrightarrow}}E$.
\end{enumerate}

B) Conditions équivalentes
\begin{enumerate}[label={},leftmargin=1.8em]
\item (i) $\overleftarrow{\alpha}_*(1_E)=1_{E'}$
\item (ii) $\unsure{\overleftarrow{\beta}_*|\widetilde{\beta}_*}(\unsure{S})=\varnothing_E\Longrightarrow S=\varnothing_{E'}$.
\item (iii) $f'$ surjectif --- i.e. $\alpha$ \unsure{composé} d'un comorphisme, et d'un morphisme surjectif.
\end{enumerate}

\uline{Corollaire 2} Conditions équivalentes
\begin{enumerate}[label={},leftmargin=1.8em]
\item (i) $\overleftarrow{\alpha}_*$ commute à $|\unsure{\circ|0|\bullet}|$ ``\unsure{fidèlement}'', i.e. $\forall S,T\in\unsure{\mathcal{P}|\mathscr{P}}(E)$,
\[
S\,|\circ|\,T\Longleftrightarrow\bar\alpha(S)\,|\circ|\,\bar\alpha(T)
\]
\item (ii) $\overleftarrow{\alpha}_*$ commute à $|\circ|$ (cf A) (ou une des conditions équivalentes à celle-ci), et
\[
\overleftarrow{\alpha}_*(S)=\varnothing\Longrightarrow S=\varnothing.
\]
\item (iii) $\alpha$ est \unsure{composé} d'un comorphisme \unsure{défini} par une application \uline{surjective} $f$, et d'un morphisme injectif.
\end{enumerate}

\uline{Corollaire 3} Conditions équivalentes
\begin{enumerate}[label={},leftmargin=1.8em]
\item (i) $\alpha_*$ commute à $C$ \inserted{(cf A) pour cond. équiv.)}, et $\alpha_*(S)=0_{E'}\Rightarrow S=0_E$
\item (ii) $\alpha_*$ commute à $|\circ|$ \unsure{fidèlement}, et $1\to1$
\item (iii) $\alpha$ est un comorphisme \unsure{correspondant} à une application \uline{surjective} $g:E'\mathrel{-}E$.
\end{enumerate}

\scanpage{107}
\mspage{103}
% archive sheet number: 106

La question se pose dans quelle mesure ces équivalences \deleted{sont valables} \unsure{gardent} un sens et restent valables pour une correspondance entre disjoncteurs généraux. Faute d'avoir une \unsure{décomposition} du type $\alpha_*=f'_*\,{}^{t}f_*$, \deleted{il y a} l'interprétation ``naïve'' ou ``\unsure{directe}'' en termes des propriétés des $f$, $f'$, \deleted{\unsure{va}} semble destinée à disparaître. En fait, ce sont les équivalents algébriques ``\unsure{rappelés}'' tantôt, qui sont en \unsure{premier lieu|première ligne}.

% margin note written sideways
% slanted
\margin{Voir \illegible{1} \unsure{$2^{\mathrm{me}}$ pt} \unsure{CF} \uline{VIII} \unsure{exemple} \unsure{non général} (\unsure{cont.} \unsure{top.}) \unsure{72} \unsure{$\Downarrow|\downarrow$} $\ldots$}
Malheureusement, je ne sais prouver les équivalences de i) à (iii') de A) \inserted{(\unsure{ou} que i) \unsure{implique} \unsure{(ii)} (iii))}, que moyennant des hypothèses draconiennes, proches déjà de l'hypothèse booléenne (\unsure{cf} \unsure{donc} \unsure{aux.} \unsure{Axiomes|Nécessités} des \unsure{repères} à propos des supports). Peut-être est-ce là encore des faux problèmes, vu que pour mes propos géométriques, on ne travaillera jamais qu'avec des supports très particuliers (essentiellement, les supports constructibles) : et on n'aura à prendre des $\operatorname{Inf}$ et $\unsure{=}$ des $\operatorname{Sup}$ que sous des conditions idoines de ``\unsure{commensurabilité}''. Ce sont les cas où, précisément, les arguments

\scanpage{108}
\mspage{104}
% archive sheet number: 107

auxquels j'ai fait allusion ont \unsure{marché}, dans A), pour nous dire que la commutation des $\alpha_*$ à $|\unsure{\circ|0|\bullet}|$ implique les autres propriétés (sous les hyp. habituelles de \unsure{commensurabilité}). Donc l'interprétation géom. attachée à ``$\unsure{u_\alpha|u_*}$ commute à $|\circ|$'' sera bien celle d'un ``comorphisme'' $f^{-1}$, suivi d'une application injective. Quand \uline{de plus} $u_\alpha(1)=1$, on s'attendra donc à avoir commutation à $\unsure{C|\complement}$ (et on pourra le prouver, moyennant des conditions de \unsure{constructibilité} \unsure{convenables})

\begin{quote}
\noindent\vrule\quad\begin{minipage}{0.88\linewidth}
NB. La commutation à $C$ implique tjrs que $1\longrightarrow1$, ainsi que la compatibilité à $|\circ|$.
\end{minipage}
\end{quote}

\uuline{Question} Si $\alpha_*$ \uline{et} $\beta_*$ commutent à $C$, est-ce des iso. de spatios, \inserted{\unsure{(oui.)|(oui.}} (inverses l'un de l'autre) ?

Notons enfin que l'équivalence de (i) et (ii) dans (B) se vérifie trivialement dans tout spatios.

\scanpage{109}
\mspage{105}
% archive sheet number: 108

\uline{Lemme 1} Soit $b\leq c$ dans $\unsure{\Sigma|\bar\Sigma}$. Conditions équivalentes
\begin{enumerate}[label={},leftmargin=1.8em]
\item (i) $\exists b^*\in\Sigma$, avec $b^*\,|\unsure{\circ|0|\bullet}|\,b$ et $b\vee b^*=c$
\item (ii) On a $b\vee b_0^*=c$, où $b_0^*=\unsure{C|\complement}b\wedge c$
\end{enumerate}
De plus, si on a $\unsure{b'|b^*}$ comme en (i), on a $b^*\leq b_0^*$, i.e. $b_0^*$ est le plus grand des $b^*$ satisfaisant (i).

\uline{Lemme 2} Soit $b\leq c$ dans $\Sigma$. Conditions équivalentes
\begin{enumerate}[label={},leftmargin=1.8em]
\item (i) $b,c$ engendrent dans $\Sigma$ \inserted{(pour $\vee,\wedge,C$)} une sous-algèbre de Boole.
\item (ii) $b,c$ engendrent une sous-algèbre de Boole finie, de cardinal $2^n$ avec $n\in\{1,2,3\}$.

\inserted{NB} (($n\deleted{\unsure{\leq}}=1$ sss $\{b,c\}\subset\{0,1\}$ ; $n\leq2$ sss $b=0$ ou $c=1$ ou $b=c$ ; \deleted{mais $\{b,c\}\ne\{0,1\}$} \inserted{\deleted{$\{b,c\}\ne\{0,1\}$ (donc exclus}} \deleted{\unsure{mais} $c\ne0$, $b\ne1$ et les cas $b=\unsure{c|b}=0$, $b=\unsure{c|b}=1$, \unsure{ou} \inserted{et} $b=0$ et $c=1$)} ; $n=3$ sss $0<b<c<1$).
\item \unsure{(iii)|(iii')} $(b,c)$ satisfont les conditions du lemme 1, et $(c',b')$ aussi (où $c'=Cc$, $b'=Cb$).
\item \inserted{(iv) Posant $\alpha=b$, $\beta=c\wedge Cb$, $\unsure{\gamma|\sigma}=Cc$, \unsure{les|b,} $\alpha,\beta,\gamma$ engendrent une sous-alg. de Boole \deleted{\illegible{1}} \unsure{ou encore, on a les} \deleted{\unsure{aux}} relations $C(\alpha\vee\beta)=\gamma$, $C(\beta\vee\gamma)=\alpha$ \inserted{$C(\gamma\vee\alpha)=\beta$} / i.e. $C(\alpha)\wedge(\gamma\vee\alpha)=\gamma$, $C(\unsure{\beta|\gamma})(\gamma\vee\alpha)=\alpha$}
% slanted
\margin{\unsure{les} $\alpha,\beta,\gamma$ sont disjoints}
\end{enumerate}

\uline{Cor.} Conditions équivalentes sur $\Sigma$
\begin{enumerate}[label={},leftmargin=1.8em]
\item (i) $\forall$ couples $b,c$ avec $b\leq c$, les conditions du lemme 1 satisfaites.
\item (ii) $\forall$ couples $b,c$ avec $b\leq c$, les conditions \inserted{(plus fortes)} du lemme 2 sont satisfaites.
\end{enumerate}
% margin note written sideways
% slanted
\margin{\deleted{\unsure{ici}} \unsure{On peut voir|On peut remplacer} une famille finie $0\leq b_1\leq b_2\mathrel{\unsure{-|\cdots}}\leq b_n\leq1$ (pour faire (ii'))}
% margin note written sideways
% slanted
\margin{\unsure{dégager} \unsure{lemme} 3 \unsure{pour compléter}, \deleted{\unsure{articuler}} (\unsure{en} Cor. 2) \unsure{plus bas}}
\noindent$\left[\rule{0pt}{10em}\right.$\begin{minipage}{0.88\linewidth}
(iii) $\forall$ couples $\alpha,\beta$ avec $\alpha\,|\circ|\,\beta$, $\alpha$ et $\beta$ engendrent une sous-algèbre de Boole (laquelle sera \unsure{automatiquement} finie de cardinal $2^n$, $\deleted{1\leq n\leq}n\in\{1,2,3\}$), $n=1$ sss $\{\alpha,\beta\}\unsure{\in|\subset}\{0,1\}$, $n\leq2$ sss $\alpha$ ou $\beta$ est nul \inserted{ou $\alpha\vee\beta=1$} \deleted{\unsure{pour tous}} \deleted{sss \unsure{(NB)} $\alpha\ne0$, \unsure{et}} \deleted{\unsure{alors}} $n=3$ sss $\alpha\ne0$, $\beta\ne0$, $\alpha\vee\beta\ne1$.

\inserted{\unsure{(iiii)|(iiii')|(iii')|(iv)} Comme (iii), mais avec familles \uline{finies} $(\alpha_i)$, avec $\alpha_i\,|\circ|\,\alpha_j$ pour $i\ne j$}
\end{minipage}

\scanpage{110}
\mspage{NA 21}
% archive sheet number: 109
% The typescript is upside down in the supplied scan; reading order below is after rotation by 180 degrees.
% A continuous pen stroke traverses the body of the reused typescript. Its cancellation is represented paragraph by paragraph; nested deletions are local typewritten cancellations.

\deleted{formulation précise en termes d'une notion de ``régularité'' ad hoc (dont Kashiwara ne disposait pas, pas en 1975,76 tout au moins), et enfin dans la démonstration\deleted{\unsure{,}} du théorème, dont les principaux ingrédients avaient \unsure{été} réunis déjà dans sa thèse (dans le cas\unsure{-| }\unsure{soeux|ceux} des \handwritten{${}^{\unsure{\infty|\mathscr{D}}}$}-Modules), à un moment où personne, apparemment, ne se souciait encore le moins du monde de gaspiller un temps précieux à vouloir prouver un énoncé ressenti comme \inserted{accidentel,} marginal.}

\deleted{Le lecteur attentif fera de même les \unsure{ajustements|ajustement} qui s'imposent, quand par la suite il m'arrivera de parler de la ``philosophie de Mebkhout'' des -Modules. Il semble bien, par les documents et témoignages à ma disposition, qu'il ait été le seul, \inserted{entre 1976 et 1980,} à se consacrer corps et âme au développement de cette nouvelle philosophie des coefficients\deleted{\unsure{,}} \deleted{\unsure{mixtes}} en train de naître cahin-caha, et à la faire connaître, et ses publications sont à ce sujet un témoignage éloquent. Cela n'est en aucun cas une raison valable, cependant, pour faire mine d'ignorer pour autant \deleted{\unsure{la}}le \deleted{\unsure{contribution}}\inserted{rôle} de Kashiwara dans l'éclosion de certaines idées importantes, même si ce rôle n'est attesté par aucune publication en bonne et due forme, mais étaient transmises simplement (comme cela a été le cas si souvent pour \unsure{moi-même} dans les années soixante) de bouche à oreille. \deleted{Si donc il \unsure{devait|venait} \illegible{2} chercher un nom à cette philosophie-là (autre que le nom bâtard de ``Riemann-Hilbert'' !), \illegible{4} \unsure{Pour ma part}, j'espère qu'il lui \unsure{restera}.}}

\deleted{Page \unsure{963|953}, deux compléments à la note de haut de page:}

\deleted{($^{x}$) (6 mai 1986) Il y a lieu, visiblement, d'exclure le nom de M. Kashiwara de la liste des participants à la ``mystification-escroquerie'' du Colloque Pervers''. S'il est vrai que ``l'échange de mauvais procédés'' entre Kashiwara et Mebkhout était alors engagé depuis belle lurette, il est vrai que non seulement Kashiwara est étranger à la mise en scène du Colloque Pervers, mais qu'il en est, à part Mebkhout\deleted{\unsure{,}} \deleted{\unsure{le principal}} (et l'ancêtre, il va sans dire), le principal dindon de la farce, vu qu'il n'y est pas plus question de lui, que de l'inconnu de service alias Mebkhout. (Il est vrai que c'était de moindre conséquence pour lui, qui depuis longtemps fait partie de l'establishment mathématique, que pour Mebkhout, pratiquement inconnu et ``seul contre tous''.)}

\deleted{($^{xx}$) (6 mai 1986) Pierre Schapira \inserted{me signale} que l'idée de la transformée de Fourier dans le cadre des -Modules avait été développée (mais pas sous ce nom) dans l'article de \unsure{Seto|Sato},Kawai, Kashiwara de \unsure{1973|1975} déjà cité (note de b. de p. ($\unsure{xxx|***}$) p. 948), donc bien avant les travaux de Malgrange sur ce thème.}

\scanpage{111}
\mspage{106}
% archive sheet number: 110

\uline{Corollaire 2} \inserted{Sous les conditions \inserted{\unsure{équ.}} du Cor.\,1 (iii),} Si $\alpha,\beta\ne0$, $\alpha\vee\beta\ne1$, et posant $\unsure{\gamma|\sigma}=\unsure{C|\complement}(\alpha\vee\beta)\ne0$, $\alpha,\beta,\gamma$ forment une ``base'' de l'algèbre de Boole engendrée par $\alpha,\beta$. Sous les conditions du lemme 2, \inserted{et supposant $0<b<c<1$,} posant
\[
\alpha=b,\quad\beta=c\wedge Cb,\quad\gamma=Cc
\]
$\alpha,\beta,\gamma$ forment une base de l'algèbre de Boole engendrée par $b,c$.

\noindent\makebox[0pt][r]{\raisebox{-4\baselineskip}[0pt][0pt]{\rule{0.3pt}{5\baselineskip}\kern2pt\rule{0.3pt}{5\baselineskip}}\hspace{0.8em}}\uline{Question} Les conditions envisagées au \deleted{lemme 1} Cor.\,1 sur $\unsure{\Sigma|\bar\Sigma}$ sont elles tjrs vérifiées ? J'en doute fort ! Le sont elles si $\Sigma$ est fini${}^{*}$ ? J'en doute également.
% slanted
\margin{Non, \unsure{contrexemple|contre-exemple} avec $\operatorname{Card}\Sigma=10$}
Si on suppose ces conditions satisfaites, $\Sigma$ est-il nécessairement une algèbre de Boole ? L'est-il si on fait l'hypothèse similaire à celle du Cor.\,1 \inserted{(\unsure{iiii'|iiii|iv})}, pour \uline{$\unsure{\forall|\text{tt}}$} familles $(\alpha_i)_{\unsure{\nu|i}\in I}$ d'éléments mutuellement disjoints de $\unsure{I|\Sigma}$ (pas néc. finie).

\begin{quote}
\noindent\vrule\kern2pt\vrule\quad\begin{minipage}{0.86\linewidth}
Il faudrait absolument avoir le coeur net de ce qui en est, quand $\Sigma$ est défini via un ens. ordonné $L$, muni de la relation $\unsure{a|x}|\unsure{\circ|0|\bullet}|y$ sss $x<y$ ou $y<x$.
\end{minipage}
\end{quote}

